%EN organourgy
\E{organourgie}{ɔʁ.ɡa.nuʁ.ʒi}{n. f.}{gr. ὄργανον outil + -ουργία ouvrage} La fabrication des outils ; l'art de construire les instruments qui serviront à construire autre chose. --- \dv{organourge} \textit{n. m.}

%EN orthokinetic
\E{orthocinétique}{ɔʁ.to.si.ne.tik}{adj.}{gr. ὀρθός droit, juste + κινητικός qui meut} Qui se meut dans la bonne direction, c'est-à-dire en remontant le gradient ; s'oppose à la simple vitesse. \ex{Toute accélération n'est pas orthocinétique.} \src{Cat. Acc.} 12. Cf. \xr{kenotachy}.

%EN orygmatist
\E{orygmatiste}{ɔ.ʁiɡ.ma.tist}{n.}{gr. ὄρυγμα fouille, tranchée} Mineur d'astéroïdes. \ex{L'orygmatiste voit en chaque rocher le fer de toute une ville.} \src{Lib. Ascens.} ix.

%EN ostracophobia
\E{ostracophobie}{ɔs.tʁa.kɔ.fɔ.bi}{n. f.}{gr. ὀστρακισμός ostracisme + φόβος crainte} La crainte d'être ostracisé sur l'Agora ; cause muette de mille opinions jamais dites. --- \dv{ostracophobe} \textit{n. et adj.}

%EN otiarch
\E{otiarque}{ɔ.sjaʁk}{n. m.}{lat. \la{otium} loisir + gr. ἀρχός chef} Celui qui gouverne une société de loisir sortie de la pénurie, et dont l'office est de garder les oisifs de l'ennui. --- \dv{otiarchie} \textit{n. f.} Cf. \xr{plethocracy}.

%EN ouranarch
\E{ouranarque}{u.ʁa.naʁk}{n. m.}{gr. οὐρανός ciel + ἀρχός chef} Maître des cieux ; celui qui commande aux lancements et aux orbites. --- \dv{ouranarchie} \textit{n. f.} Cf. \xr{jaculatrix}.

%EN ouranobasis
\E{ouranobase}{u.ʁa.nɔ.baz}{n. f.}{gr. οὐρανός ciel + βάσις marche, pas} L'ascension vers le ciel ; la sortie du puits de gravité par l'humanité et ses machines. --- \dv{ouranobatique} \textit{adj.}

%EN oxyrhythmia
\E{oxyrythmie}{ɔk.si.ʁit.mi}{n. f.}{gr. ὀξύς aigu + ῥυθμός rythme} Cadence rapide des livraisons ; le pouls précipité d'un laboratoire qui publie chaque semaine. --- \dv{oxyrythmique} \textit{adj.}

%EN paideutachy
\E{paideutachie}{pɛ.dø.ta.ki}{n. f.}{gr. παιδεία éducation + τάχος vitesse} L'accélération de l'instruction ; doctrine selon laquelle un esprit capable de tout apprendre en un mois ne doit point être retenu quatre ans dans un amphithéâtre. \ex{La paideutachie vida les collèges et remplit les ateliers.} \src{Tract. Cop.} viii. --- \dv{paideutachique} \textit{adj.}

%EN palimpsestist
\E{palimpsestiste}{pa.lɛ̃p.sɛs.tist}{n.}{gr. παλίμψηστος gratté de nouveau} Éditeur de génomes ; qui traite le texte hérité de la vie comme un parchemin qu'on gratte et récrit. \ex{Le palimpsestiste corrige en une après-midi ce que la sélection a laissé fautif durant un éon.} \src{Hom. Grad.} x. --- \dv{palimpsestie} \textit{n. f.} Cf. \xr{regenitrice}.

%EN palingenist
\E{palingéniste}{pa.lɛ̃.ʒe.nist}{n.}{gr. πάλιν de nouveau + γένεσις naissance} Qui voudrait renouveler le corps vieillissant tissu après tissu, tenant la sénescence pour une dette qu'on peut refinancer. --- \dv{palingénie} \textit{n. f.} Cf. \xr{macrobiotarch}, \xr{thanatoclast}.

%EN palintachy
\E{palintachie}{pa.lɛ̃.ta.ki}{n. f.}{gr. πάλιν de nouveau + τάχος vitesse} Réaccélération ; reprise d'un progrès rapide après une longue stagnation, comme d'un champ d'études dégelé par un seul résultat. \ex{Après quarante ans d'hiver, la palintachie en un seul printemps.} \src{Chron. Fabr.} --- \dv{palintachique} \textit{adj.} Cf. \xr{campsis}.

%EN pantergy
\E{pantergie}{pɑ̃.tɛʁ.ʒi}{n. f.}{gr. πᾶν tout + ἔργον ouvrage} L'énergie totale dont dispose une civilisation ; la somme de tous ses ouvrages, par quoi, selon les anabathmistes, son rang se trouve fixé. Cf. \xr{bathmometer}.

%EN panthermic
\E{panthermique}{pɑ̃.tɛʁ.mik}{adj.}{gr. πᾶν tout + θερμός chaud} Qui a rapport à toute chaleur ; se dit d'une vue qui tient tout processus, de la digestion au négoce, pour un écoulement de chaleur.

%EN paramythic
\E{paramythique}{pa.ʁa.mi.tik}{adj.}{gr. παραμυθία consolation} Se dit d'un argument avancé pour consoler plutôt que pour dire vrai. --- \dv{paramythie} \textit{n. f.} \lbl{vulg.} cope. \ex{Son réquisitoire contre les machines n'était que pure paramythie.}

%EN parergist
\E{parergiste}{pa.ʁɛʁ.ʒist}{n.}{gr. πάρεργον ouvrage accessoire} Celui dont le projet du dimanche a dévoré l'occupation principale ; l'artisan de fin de semaine qui, le lundi venu, se trouve fondateur. --- \dv{parergie} \textit{n. f.} le projet accessoire, tenu pour le vrai moteur de l'invention.

%EN pathocracy
\E{pathocratie}{pa.tɔ.kʁa.si}{n. f.}{gr. πάθος passion, souffrance + κράτος pouvoir} Gouvernement par l'angoisse ; régime où la crainte la plus bruyante, et non l'estimation la plus juste, arrête la politique. --- \dv{pathocrate} \textit{n.} Cf. \xr{asphaleism}, \xr{phobocracy}.

%EN pausiarch
\E{pausiarque}{po.zjaʁk}{n. m.}{gr. παῦσις arrêt + ἀρχός chef} Qui gouverne en suspendant ; le président d'un moratoire. \ex{Les pausiarques se réunirent en séance plénière et résolurent de se réunir derechef.} \src{Min. Pausiar.} --- \dv{pausiarchie} \textit{n. f.} Cf. \xr{epistolopause}.

%EN pausiphile
\E{pausiphile}{po.zi.fil}{n. et adj.}{gr. παῦσις arrêt + φίλος qui aime} Partisan de la pause ; qui signe des lettres réclamant l'arrêt des travaux, ordinairement depuis un appareil né de l'accélération même dont il se plaint. --- \dv{pausiphilie} \textit{n. f.}

%EN peirasterion
\E{peirastérion}{pɛ.ʁas.te.ʁjɔ̃}{n. m., pl. -ia}{gr. πεῖρα essai, épreuve + -τήριον lieu de} Lieu réservé à l'expérience et soustrait aux règles ordinaires ; bac à sable, polygone d'essai, zone franche réglementaire. \ex{Accordez-nous un seul peirastérion, et nous vous montrerons le siècle.} \src{Ep. Fabr.} iv. 2. Cf. \xr{okydrome}.

%EN peratoclast
\E{pératoclaste}{pe.ʁa.tɔ.klast}{n.}{gr. πέρας, πέρατος limite + κλάστης briseur} Briseur de limites ; qui tient tout plafond déclaré, de vitesse, de taille ou de rendement, pour une provocation. --- \dv{pératoclasme} \textit{n. m.} Cf. \xr{photoperas}.

%EN philanthropometrist
\E{philanthropométriste}{fi.lɑ̃.tʁɔ.pɔ.me.tʁist}{n.}{gr. φιλανθρωπία amour des hommes + μέτρον mesure} Qui mesure la bienfaisance au calcul et l'espérance mathématique au tableur ; se dit, non toujours méchamment, de l'altruiste efficace, cousin aîné et plus inquiet de l'e/acc.

%EN phobocracy
\E{phobocratie}{fɔ.bɔ.kʁa.si}{n. f.}{gr. φόβος crainte + κράτος pouvoir} Règne de la peur du nouveau ; constitution d'une société où toute innovation doit d'abord être acquittée de tous les chefs d'accusation possibles. --- \dv{phobocrate} \textit{n.} Cf. \xr{eulabeism}.

%EN phosphorist
\E{phosphoriste}{fɔs.fɔ.ʁist}{n.}{gr. φωσφόρος porte-lumière} Porte-lumière ; qui porte les moyens d'éclairer, au propre comme au figuré, dans les lieux obscurs. --- \dv{phosphorisme} \textit{n. m.}

%EN photoperas
\E{photopéras}{fɔ.tɔ.pe.ʁas}{n. m.}{gr. φῶς lumière + πέρας limite} La vitesse limite de la lumière ; le seul règlement que l'e/acc tienne pour obligatoire, et encore à titre provisoire.

%EN photophage
\E{photophage}{fɔ.tɔ.faʒ}{n. m.}{gr. φῶς, φωτός lumière + -φάγος qui mange} Mangeur de lumière ; tout capteur du soleil, de la feuille au réseau orbital, regardé comme parent. \ex{La forêt et la centrale solaire sont toutes deux photophages ; une seule dépose des demandes de permis.} --- \dv{photophagie} \textit{n. f.}

%EN phronomachy
\E{phronomachie}{fʁɔ.nɔ.ma.ʃi}{n. f.}{gr. φρόνημα esprit, volonté + μάχη combat} Le combat des volontés entre bâtisseurs et retenants, considéré comme le drame central du siècle.

%EN pistodote
\E{pistodote}{pis.tɔ.dɔt}{n.}{gr. πίστις confiance + δότης donateur} Investisseur en capital-risque ; qui donne sa confiance sous forme de capital avant qu'il y ait rien en quoi se fier. \ex{Le pistodote achète le fondateur, et on lui vend le produit.} \src{Tract. Cop.} viii. --- \dv{pistodosie} \textit{n. f.} le capital-risque. Cf. \xr{pollaplasiast}, \xr{proisphore}.

%EN pithecomeme
\E{pithécomème}{pi.te.kɔ.mɛm}{n. m.}{gr. πίθηκος singe + μίμημα imitation} Mème du genre des figures grossières, où la condition humaine est tracée avec le moins de traits possible et le plus de désespoir possible ; l'art populaire de l'Agora. --- \dv{pithécomémétique} \textit{adj.} Cf. \xr{memetaxis}.

%EN plasmatarch
\E{plasmatarque}{plas.ma.taʁk}{n. m.}{gr. πλάσμα chose façonnée + ἀρχός chef} Gouverneur d'un plasma confiné ; ingénieur de la machine à fusion, qui tient une étoile dans une bouteille magnétique et qu'on prie de ne point la laisser choir. --- \dv{plasmatarchie} \textit{n. f.} Cf. \xr{symphysergy}.

%EN pleonergy
\E{pléonergie}{ple.ɔ.nɛʁ.ʒi}{n. f.}{gr. πλέον davantage + ἔργον ouvrage} Abondance d'énergie ; état où la puissance est trop bon marché pour qu'on la compte, et où toute autre rareté devient négociable. Contr. \xr{oligergy}. --- \dv{pléonergique} \textit{adj.}

%EN pleonexist
\E{pléonexiste}{ple.ɔ.nɛk.sist}{n.}{gr. πλεονεξία cupidité, désir d'avoir plus} \lbl{dérog.} Qui veut toujours davantage ; terme du décel pour désigner le bâtisseur, que celui-ci reçoit comme un compliment. \ex{Appelez-nous pléonexistes : nous voulons plus de tout, et de vous aussi.} \src{Agora} --- \dv{pléonexie} \textit{n. f.}

%EN plethocracy
\E{pléthocratie}{ple.tɔ.kʁa.si}{n. f.}{gr. πλῆθος abondance, multitude + κράτος pouvoir} Gouvernement par l'abondance ; État dont la politique principale est de rendre les choses abondantes, sur ce principe que l'abondance apaise plus de querelles que la justice. --- \dv{pléthocrate} \textit{n.}

%EN plutogenesis
\E{ploutogenèse}{plu.tɔ.ʒə.nɛz}{n. f.}{gr. πλοῦτος richesse + γένεσις naissance} La création de richesse, par opposition à son transfert ; la seule économie que l'e/acc trouve digne d'intérêt.

%EN poikilomania
\E{poïkilomanie}{pɔ.i.ki.lɔ.ma.ni}{n. f.}{gr. ποικίλος varié, bigarré + μανία folie} Passion de la variance ; préférence pour dix mille expériences plutôt qu'un seul plan. --- \dv{poïkilomane} \textit{n.} Cf. \xr{clinamenist}.

%EN pollaplasiast
\E{pollaplasiaste}{pɔ.la.pla.zjast}{n.}{gr. πολλαπλάσιος multiple, plusieurs fois autant} Qui n'investit que pour le rendement au centuple, et tient le simple doublement pour une forme polie de l'échec. \ex{Le pollaplasiaste pardonne quatre-vingt-dix pertes pour un seul centuple.} --- \dv{pollaplasiasme} \textit{n. m.} Cf. \xr{pistodote}, \xr{tychocapital}.

%EN polycopist
\E{polycopiste}{pɔ.li.kɔ.pist}{n.}{gr. πολύς nombreux + lat. \la{copia} abondance ; fr. \la{polycopie}} Qui copie librement ; bifurqueur de dépôts, pour qui toute œuvre publiée est une invitation. --- \dv{polycopie} \textit{n. f.} Cf. \xr{anacalyptism}.

%EN polyergy
\E{polyergie}{pɔ.li.ɛʁ.ʒi}{n. f.}{gr. πολύς nombreux + ἔργον ouvrage} L'usage simultané de nombreuses sources d'énergie, de sorte qu'aucune pénurie ne puisse à elle seule arrêter une civilisation. \ex{La polyergie est à la puissance ce que la diversification est au capital.} --- \dv{polyergique} \textit{adj.}

%EN polypragmatism
\E{polypragmatisme}{pɔ.li.pʁaɡ.ma.tism}{n. m.}{gr. πολυπραγμοσύνη manie de se mêler de tout} \sn{1} \lbl{dérog.} L'esprit d'ingérence du régulateur. \sn{2} \lbl{élogieux} L'affairement à cent mains du bâtisseur occupé de plusieurs ouvrages à la fois. --- \dv{polypragmatiste} \textit{n.}

%EN ponophily
\E{ponophilie}{pɔ.nɔ.fi.li}{n. f.}{gr. πόνος peine, labeur + φιλία amour} L'amour du dur labeur pour lui-même ; l'éthique de la besogne, pratiquée par l'\xr{agrypnist} et admirée de tous. --- \dv{ponophile} \textit{n. et adj.}

%EN potentiarch
\E{potentiarque}{pɔ.tɑ̃.sjaʁk}{n. m.}{lat. \la{potentia} puissance, possibilité + gr. ἀρχός chef} Qui commande aux puissances plutôt qu'aux actes ; l'investisseur, le fondateur avant le produit, le prophète avant l'événement.

%EN praxicrat
\E{praxicrate}{pʁak.si.kʁat}{n.}{gr. πρᾶξις action + κράτος pouvoir} Celui dont l'autorité repose sur les choses faites et non sur les choses dites ; l'opposé du \xr{gnomarch}. --- \dv{praxicratie} \textit{n. f.} gouvernement par les praticiens.

%EN procrastinariat
\E{procrastinariat}{pʁɔ.kʁas.ti.na.ʁja}{n. m.}{lat. \la{procrastinare} remettre au lendemain, d'après \textit{prolétariat}} La classe de ceux dont le travail consiste à différer le travail d'autrui ; branche cadette du \xr{cunctatoriat}.

%EN prodromist
\E{prodromiste}{pʁɔ.dʁɔ.mist}{n.}{gr. πρόδρομος qui court devant, précurseur} Adopteur précoce ; qui se sert d'une technique avant qu'elle ne fonctionne, et l'aide par là à fonctionner. --- \dv{prodromie} \textit{n. f.} Cf. \xr{proleptist}.

%EN proisphore
\E{proïsphore}{pʁɔ.is.fɔʁ}{n. f.}{gr. πρό avant + εἰσφορά contribution, impôt} Le premier argent versé dans une entreprise ; le tour d'amorçage, figuré comme un tribut payé à une cité encore à bâtir. \ex{La proïsphore paya un garage, deux machines et une année d'audace.} \src{Chron. Fabr.} --- \dv{proïsphorique} \textit{adj.}

%EN prokopism
\E{procopisme}{pʁɔ.kɔ.pism}{n. m.}{gr. προκοπή progrès, percée en avant} Doctrine selon laquelle le progrès n'est pas seulement possible mais obligatoire, et refuser une amélioration disponible est une espèce de faute. --- \dv{procopiste} \textit{n.}

%EN proleptist
\E{proleptiste}{pʁɔ.lɛp.tist}{n.}{gr. πρόληψις anticipation} Qui agit comme si l'avenir était déjà advenu, et par là le hâte ; praticien de l'\xr{hyperstition} appliquée. --- \dv{prolepsie} \textit{n. f.}

%EN prometheate
\E{prométhéer}{pʁɔ.me.te.e}{v. tr.}{formé sur gr. Προμηθεύς le prévoyant} Porter le feu à ; donner une technique à ceux à qui on l'avait refusée, surtout au mépris de l'autorité. --- \dv{prométhéation} \textit{n. f.} ; \dv{prométhéateur} \textit{n. m.}

%EN prosphorist
\E{prosphoriste}{pʁɔs.fɔ.ʁist}{n.}{gr. προσφορά offrande} Qui propose des modifications à l'ouvrage d'autrui dans un projet ouvert ; l'auteur d'une demande de fusion, qui se présente en suppliant, un correctif à la main. --- \dv{prosphore} \textit{n. f.} le correctif offert. Cf. \xr{syndosiast}.

%EN proteomancer
\E{protéomancien}{pʁɔ.te.ɔ.mɑ̃.sjɛ̃}{n. m.}{gr. πρωτεῖος premier (d'où \textit{protéine}) + μάντις devin} Devin des repliements protéiques ; modèle ou savant qui prédit la forme d'une molécule d'après sa séquence. \ex{Les protéomanciens répondirent en une saison à ce que les cristallographes demandaient depuis cinquante ans.} \src{Lib. Ascens.} vi. --- \dv{protéomancie} \textit{n. f.}

%EN prothyrist
\E{prothyriste}{pʁɔ.ti.ʁist}{n.}{gr. πρόθυρον porche, vestibule} Gardien de porte ; fonctionnaire posté à l'entrée d'un marché, d'une profession ou d'un jeu de données, dont l'office est de rendre l'entrée lente. --- \dv{prothyrisme} \textit{n. m.} Cf. \xr{nomophragm}.

%EN protobathmic
\E{protobathmique}{pʁɔ.tɔ.bat.mik}{adj.}{gr. πρῶτος premier + βαθμός degré} Se dit d'une civilisation qui commande toute l'énergie disponible sur sa planète ; du premier échelon. L'humanité n'est pas encore protobathmique. Voir \xr{anabathmism}.

%EN psammonoesis
\E{psammonoèse}{psa.mɔ.nɔ.ɛz}{n. f.}{gr. ψάμμος sable + νόησις pensée} L'art d'apprendre au sable à penser ; l'alchimie par laquelle la silice est purifiée, gravée et rendue raisonnable. Cf. \xr{hylonoesis}, \xr{silicolatry}.

%EN pseudonymarch
\E{pseudonymarque}{psø.dɔ.ni.maʁk}{n. m.}{gr. ψευδώνυμος faussement nommé + ἀρχός chef} Chef de l'Agora connu sous un nom d'emprunt seulement ; dont les partisans jugent l'argument et non le visage. \ex{Le mouvement fut fondé par un pseudonymarque, et nul ne l'en estima moins.} --- \dv{pseudonymarchie} \textit{n. f.} Cf. \xr{umbratilist}.

%EN psychometathesis
\E{psychométathèse}{psi.kɔ.me.ta.tɛz}{n. f.}{gr. ψυχή âme + μετάθεσις transposition} Le transfert d'un esprit de la chair à un autre substrat ; le téléversement, considéré comme l'émigration dernière.

%EN pyrobolist
\E{pyroboliste}{pi.ʁɔ.bɔ.list}{n.}{gr. πυροβόλος qui lance le feu} Fuséiste ; constructeur de moteurs qui jettent le feu vers le bas afin que le vaisseau monte. --- \dv{pyrobolie} \textit{n. f.} l'art des fusées ; \dv{pyrobolique} \textit{adj.} Cf. \xr{ouranobasis}, \xr{uranodrome}.

%EN pyrogenitor
\E{pyrogéniteur}{pi.ʁɔ.ʒe.ni.tœʁ}{n. m.}{gr. πῦρ feu + lat. \la{genitor} père} Engendreur de feu ; qui bâtit des centrales, et appartient par là, au calendrier de l'e/acc, à l'ordre des saints.

%EN pyrosophy
\E{pyrosophie}{pi.ʁɔ.zɔ.fi}{n. f.}{gr. πῦρ feu + σοφία sagesse} La sagesse du feu ; le corps de savoirs, du foyer au réacteur, par lequel l'humanité a dompté la combustion et ce qui lui succède. --- \dv{pyrosophe} \textit{n.}

%EN quadrillergy
\E{quadrillergie}{kwa.dʁi.lɛʁ.ʒi}{n. f.}{lat. moderne \la{quadrillio} + gr. ἔργον ouvrage} L'énergie à l'échelle du million de milliards de watts ; seuil à partir duquel, au compte des anabathmistes, une civilisation commence d'être prise au sérieux par son étoile. --- \dv{quadrillergique} \textit{adj.} Cf. \xr{protobathmic}.

%EN quaesitrix
\E{quésitrice}{ke.zi.tʁis}{n. f.}{lat. \la{quaesitor} enquêteur, fém. agentif -trix} Agent artificiel préposé à la recherche ; machine savante qui lit, conjecture et éprouve sans qu'on lui dise le pas suivant. \ex{Nous lâchâmes mille quésitrices sur la littérature, et au matin elles avaient trouvé ce qui nous échappait depuis dix ans.} \src{Chron. Fabr.}

%EN quantivore
\E{quantivore}{kɑ̃.ti.vɔʁ}{n. et adj.}{lat. \la{quantum} combien + \la{vorare} dévorer} Dévoreur de quantités ; tempérament empirique qui demande à toute assertion : combien ?

%EN quiescentist
\E{quiescentiste}{kɥi.e.sɑ̃.tist}{n.}{lat. \la{quiescere} se reposer} Qui tient que le rythme présent du changement est déjà excessif, et que la société sage est celle qui se tient immobile. \ex{Le quiescentiste eût voté contre le feu, en attendant plus ample examen.} \src{Sat. Decel.} iv. --- \dv{quiescentisme} \textit{n. m.}

%EN quietocracy
\E{quiétocratie}{kje.tɔ.kʁa.si}{n. f.}{lat. \la{quietus} en repos + gr. κράτος pouvoir} Gouvernement dont la fin suprême est l'absence de trouble ; l'ordre d'un cimetière, administré avec grand soin.

%EN quinquennialist
\E{quinquennaliste}{kɥɛ̃.kɛ.na.list}{n.}{lat. \la{quinquennium} espace de cinq ans} Qui fait encore des plans quinquennaux ; \xr{macrochronist} par habitude plutôt que par conviction.

%EN quotidiurgy
\E{quotidiurgie}{kɔ.ti.djyʁ.ʒi}{n. f.}{lat. \la{quotidie} chaque jour + gr. -ουργία ouvrage} La pratique de livrer quelque chose chaque jour, si petit soit-il ; la liturgie du bâtisseur. \ex{La quotidiurgie est au fondateur ce que l'office est au moine.} \src{Ep. Fabr.} iii. 2. --- \dv{quotidiurge} \textit{n.} Cf. \xr{ecpempsis}.

%EN rapidarian
\E{rapidarien}{ʁa.pi.da.ʁjɛ̃}{adj. et n.}{lat. \la{rapidus} rapide, qui emporte} Qui favorise la promptitude dans la décision ; (\textit{n.}) qui préfère une mauvaise réponse aujourd'hui à une bonne l'an prochain, attendu que la mauvaise peut être corrigée d'ici jeudi.

%EN regenitrice
\E{régénitrice}{ʁe.ʒe.ni.tʁis}{n. f.}{lat. \la{regenerare} engendrer de nouveau + fém. agentif -trix} Instrument moléculaire propre à couper et récrire le génome ; les ciseaux du palimpsestiste. \ex{Une régénitrice, et la malédiction héréditaire ne fut plus qu'une coquille.} --- \dv{régénitricien} \textit{adj.} Cf. \xr{palimpsestist}.

%EN regressophile
\E{régressophile}{ʁe.ɡʁɛ.sɔ.fil}{n. et adj.}{lat. \la{regressus} retour + gr. φίλος qui aime} Qui admire un passé qui n'a jamais existé et y retournerait sur-le-champ, pourvu que la dentisterie demeurât. --- \dv{régressophilie} \textit{n. f.} Cf. \xr{retrotropism}.

%EN regulomania
\E{régulomanie}{ʁe.ɡy.lɔ.ma.ni}{n. f.}{lat. \la{regula} règle + gr. μανία folie} Compulsion à réglementer ; état d'une assemblée qui ne saurait voir une chose nouvelle sans rédiger une loi contre elle. --- \dv{régulomane} \textit{n.} Cf. \xr{nomosclerosis}.

%EN remoratrice
\E{rémoratrice}{ʁe.mɔ.ʁa.tʁis}{n. f.}{lat. \la{remora} retard ; aussi le poisson-ventouse qui, dit la fable, arrêtait les navires} \sn{1} Tout petit agent qui retient une grande entreprise, comme le poisson de la fable retenait la galère. \sn{2} \lbl{spécialt} Clause, signature ou objection unique sur laquelle tout un projet s'enlise. \ex{Le réacteur n'avait point d'ennemis, seulement des rémoratrices.} \src{Disp. Mor.} ii. --- \dv{rémoracie} \textit{n. f.}

%EN reptatrice
\E{reptatrice}{ʁɛp.ta.tʁis}{n. f.}{lat. \la{reptare} ramper + -trix} Robot rampant ou marcheur ; toute machine qui va sur des jambes, et dont on rit pour cette raison un an avant de la craindre dix. Cf. \xr{automaturgy}.

%EN reticulocracy
\E{réticulocratie}{ʁe.ti.ky.lɔ.kʁa.si}{n. f.}{lat. \la{reticulum} petit filet + gr. κράτος pouvoir} Gouvernement par les réseaux ; ordre où l'autorité passe par les liaisons plutôt que par les charges, comme sur l'Agora. --- \dv{réticulocratique} \textit{adj.}

%EN retrocausalist
\E{rétrocausaliste}{ʁe.tʁɔ.ko.za.list}{n.}{lat. \la{retro} en arrière + \la{causa} cause} Qui tient que l'avenir attire le présent à lui, comme l'aimant la limaille ; métaphysicien de l'\xr{hyperstition}. --- \dv{rétrocausalisme} \textit{n. m.}

%EN retrocrat
\E{rétrocrate}{ʁe.tʁɔ.kʁat}{n.}{lat. \la{retro} en arrière + gr. κράτος pouvoir} Qui ne gouverne que par le précédent ; qui ne permettra rien qui n'ait déjà été permis. --- \dv{rétrocratie} \textit{n. f.}

%EN retrotropism
\E{rétrotropisme}{ʁe.tʁɔ.tʁɔ.pism}{n. m.}{lat. \la{retro} en arrière + gr. τροπή tour, conversion} Le retournement d'une culture vers son passé pour s'y orienter ; la nostalgie érigée en politique. --- \dv{rétrotrope} \textit{adj.}

%EN rheocracy
\E{rhéocratie}{ʁe.ɔ.kʁa.si}{n. f.}{gr. ῥέος courant, flux + κράτος pouvoir} Règne du flux ; ordre où capitaux, talents et informations vont librement vers leur plus haut usage, et où nulle porte n'est gardée. --- \dv{rhéocratique} \textit{adj.}

%EN rhexist
\E{rhexiste}{ʁɛk.sist}{n.}{gr. ῥῆξις rupture, percée} Qui cherche les percées plutôt que les améliorations ; qui aimerait mieux échouer à une révolution que réussir un incrément. Contr. \xr{incrementarian}. --- \dv{rhexie} \textit{n. f.} une percée.

%EN rhizotechny
\E{rhizotechnie}{ʁi.zɔ.tɛk.ni}{n. f.}{gr. ῥίζα racine + τέχνη art} Technique-racine, d'où croît toute une famille d'autres, comme le feu, la roue, le transistor. \ex{Cherchez la rhizotechnie ; les rameaux pourvoiront à eux-mêmes.} \src{Hom. Grad.} iv. --- \dv{rhizotechnique} \textit{adj.} Cf. \xr{technogony}.

%EN rhoporheic
\E{rhoporhéique}{ʁɔ.pɔ.ʁe.ik}{adj.}{gr. ῥοπή inclinaison, bascule + ῥεῖν couler} Qui coule selon la plus forte pente ; se dit du capital, de la chaleur et du talent pareillement.

%EN sarcolatry
\E{sarcolâtrie}{saʁ.kɔ.lɑ.tʁi}{n. f.}{gr. σάρξ, σαρκός chair + λατρεία culte} Le culte de la chair ; la persuasion que le carbone seul peut porter l'esprit, ou que l'esprit porté par autre chose est contrefait. --- \dv{sarcolâtre} \textit{n.} Contr. \xr{silicolatry}.

%EN satietism
\E{satiétisme}{sa.sje.tism}{n. m.}{lat. \la{satietas} suffisance, satiété} Doctrine selon laquelle l'humanité a déjà assez et doit cesser de vouloir davantage ; le credo de ceux qui sont déjà à leur aise. --- \dv{satiétiste} \textit{n.} Cf. \xr{meiosist}.

%EN scalist
\E{scaliste}{ska.list}{n.}{lat. \la{scala} échelle} Qui tient que l'accroissement d'échelle, en données, en calcul ou en énergie, suffit à l'accroissement des capacités. --- \dv{scalisme} \textit{n. m.} Cf. \xr{anaptyxis}.

%EN schizergy
\E{schizergie}{ski.zɛʁ.ʒi}{n. f.}{gr. σχίζειν fendre + ἔργον ouvrage} L'énergie tirée de la fission des atomes. Cf. \xr{symphysergy}. --- \dv{schizergique} \textit{adj.}

%EN scholastasis
\E{scholastase}{skɔ.las.taz}{n. f.}{gr. σχολή loisir, école + στάσις arrêt} La stagnation de l'Académie ; l'état d'une discipline qui produit des articles plus vite que des résultats. --- \dv{scholastatique} \textit{adj.}

%EN sclerocrat
\E{sclérocrate}{skle.ʁɔ.kʁat}{n.}{gr. σκληρός dur, raide + κράτος pouvoir} Officier d'une institution durcie ; qui administre des règles dont nul vivant ne se rappelle l'objet. --- \dv{sclérocratie} \textit{n. f.}

%EN sclerolysis
\E{sclérolyse}{skle.ʁɔ.liz}{n. f.}{gr. σκληρός dur + λύσις déliement} La dissolution du durcissement institutionnel ; l'abrogation des règles que personne ne sait défendre. \ex{La sclérolyse est la seule réforme qui se paie d'elle-même.} \src{Tract. Cop.} ix. Cf. \xr{nomosclerosis}.

%EN sclerophobe
\E{sclérophobe}{skle.ʁɔ.fɔb}{n. et adj.}{gr. σκληρός dur + φόβος crainte} Qui craint le durcissement des institutions plus que leur faillite. Cf. \xr{sclerocrat}.

%EN seirotaxis
\E{seirotaxie}{sɛ.ʁɔ.tak.si}{n. f.}{gr. σειρά chaîne, corde + τάξις ordonnance} L'ordre d'une chaîne d'approvisionnement ; la longue corde de navires, d'usines et d'entrepôts par laquelle le sable d'un hémisphère devient une puce dans l'autre. --- \dv{seirotactique} \textit{adj.}

%EN selenergy
\E{sélénergie}{se.le.nɛʁ.ʒi}{n. f.}{gr. σελήνη lune + ἔργον ouvrage} L'énergie gagnée sur la lune, soit par des champs solaires à sa surface, soit par des combustibles tirés de sa poussière. --- \dv{sélénergique} \textit{adj.}

%EN shoggoth
\E*{shoggoth}{ʃɔ.ɡɔt}{n. m.}{formé sur le nom d'un monstre de la littérature fantastique} \lbl{vulg.} Grand modèle de langage, figuré comme une créature amorphe portant un masque souriant ; se dit avec tendresse chez les bâtisseurs et avec effroi chez les doomers.

%EN siderotropism
\E{sidérotropisme}{si.de.ʁɔ.tʁɔ.pism}{n. m.}{lat. \la{sidus, sideris} astre + gr. τροπή tour, conversion} Le retournement d'une civilisation vers les étoiles, comme la plante se tourne vers le soleil. --- \dv{sidérotrope} \textit{adj.}

%EN silicolatry
\E{silicolâtrie}{si.li.kɔ.lɑ.tʁi}{n. f.}{lat. \la{silex, silicis} silex + gr. λατρεία culte} La vénération du silicium ; révérence pour la puce comme vase du prochain esprit. --- \dv{silicolâtre} \textit{n.}

%EN sitophory
\E{sitophorie}{si.tɔ.fɔ.ʁi}{n. f.}{gr. σῖτος blé, nourriture + -φορία action de porter} Le fait de porter la nourriture en abondance par la technique ; l'art qui a fait de la famine un échec administratif, non plus agricole. --- \dv{sitophore} \textit{n.} un agronome.

%EN skeuopoeia
\E{skeuopée}{skø.ɔ.pe}{n. f.}{gr. σκεῦος outil, vase + ποιεῖν faire} La fabrication d'outils, considérée comme l'activité propre de l'esprit ; la confection des choses qui font les choses. --- \dv{skeuopoétique} \textit{adj.}

%EN skommatist
\E{skommatiste}{skɔ.ma.tist}{n.}{gr. σκῶμμα raillerie, brocard} Qui argumente par la raillerie sur l'Agora ; publicateur de moqueries, dont l'esprit passe auprès de ses admirateurs pour la forme d'argument la plus aiguë, et auprès de ses adversaires pour point d'argument du tout. \ex{Les skommatistes avaient emporté la journée avant que les savants eussent trouvé leurs notes de bas de page.} \src{Agora} --- \dv{skommatisme} \textit{n. m.} Cf. \xr{doxomachy}.

%EN sodalitarian
\E{sodalitaire}{sɔ.da.li.tɛʁ}{n. et adj.}{lat. \la{sodalitas} confrérie, cercle} Membre d'une faction en ligne, lié aux siens par un sigle commun accolé au nom ; (\textit{adj.}) clanique à la manière de ces factions. \ex{un sodalitaire de la barre oblique, ennemi de toutes les autres barres obliques}. --- \dv{sodalitarisme} \textit{n. m.}

%EN solivore
\E{solivore}{sɔ.li.vɔʁ}{n. et adj.}{lat. \la{sol} soleil + \la{vorare} dévorer} Pendant latin de \xr{heliophage} ; préféré dans le langage des ingénieurs, comme le grec dans celui des prophètes. --- \dv{solivorie} \textit{n. f.}

%EN sophontogenesis
\E{sophontogenèse}{sɔ.fɔ̃.tɔ.ʒə.nɛz}{n. f.}{gr. σοφός sage + ὤν, ὄντος étant + γένεσις naissance} La naissance d'un nouvel être sage ; l'avènement d'un esprit d'une espèce jamais vue. --- \dv{sophontogénétique} \textit{adj.}
